




\section{Operator Encoder-Decoder Transformer}
\label{app:architecture}

% ------------- copy from here -------------
\newcommand{\btau}{\boldsymbol{\tau}}
\newcommand{\bY}{\mathbf{Y}}
\newcommand{\bz}{\boldsymbol{z}}

In the following we describe in detail the model architecture. First we introduce the Continuum operator attention underlying our transformer architecture, then we describe the model's encoder and decoder.

% ---------------------------------------------------------------------
\subsection{Continuum Operator Attention}
\label{app:op_attn}

The standard attention mechanism implicitly treats all key positions as
equally spaced, effectively applying a uniform Riemann-sum approximation to an
underlying integral operator. On an irregular time grid this is a biased estimator.
Following \citet{calvello2024continuum}, we replace uniform averaging with a trapezoidal quadrature correction, turning each attention module into a consistent discretization of a continuum integral operator.

\paragraph{Quadrature weights.}
Given a sorted observation time grid $\tau_1 \le \tau_2 \le \cdots \le \tau_M$ of a subject (either in the context study or the target, depending on the module), define
increments $\Delta\tau_k = \tau_k - \tau_{k-1}$ with $\Delta\tau_1 = 0$,
and the trapezoidal weights
\begin{equation}
    w_k = \tfrac{1}{2}\left(\Delta\tau_k + \Delta\tau_{k+1}\right),
    \quad k = 1, \ldots, M,
    \label{eq:trap_weights}
\end{equation}
setting $\Delta\tau_{M+1} = 0$.
Padding positions are assigned weight zero and the remaining weights are
normalized to unit sum before being passed to the attention modules.

\paragraph{Operator attention.}
Let $\bm{Q} \in \mathbb{R}^{N_q \times d_A}$,
$\bm{K} \in \mathbb{R}^{M \times d_A}$,
$\bm{V} \in \mathbb{R}^{M \times d_A}$ be the query, key, and value matrices
for a single attention head, where $d_A = d / H$ is the per-head dimension, with $H$ the number of  and $d$ is the hidden dimension. The scaled dot-product attention is defined as 
\begin{equation}
{\bm S} = {\bm Q \bm K}^\top / \sqrt{d_A} + \bm{M}
\end{equation}
where $\bm{M} \in \mathbb{R}^{N_q\times M}$ is the attention masking matrix. The operator attention output at query position $i$ is given by:
\begin{equation}
    \Bigl[\mathtt{OpAttn}(\bm{Q}, \bm{K}, \bm{V}; \bm{M}, \bm{\Delta\tau})\Bigr]_i
   \ =\
    \frac{%
        \displaystyle\sum_{k=1}^{M}
        e^{S_{ik}} \cdot w_k \cdot {V}_k
    }{%
        \displaystyle\frac{1}{2}
        \sum_{k=2}^{M}
        \bigl(e^{S_{ik}} + e^{S_{i,k-1}}\bigr)\,\Delta\tau_k
    },
    \label{eq:op_attn}
\end{equation}
where the numerator is a midpoint-rule approximation of the operator 
$\int e^{q(\tau)\cdot k(\tau)}\,v(\tau)\,d\tau$
and the denominator is a trapezoidal approximation of the partition function
$\int e^{q(\tau)\cdot k(\tau)}\,d\tau$.
As the grid spacing vanishes uniformly, Eq.~\eqref{eq:op_attn} converges to
the normalized integral operator of \citet{calvello2024continuum}.
In the limit of a uniform grid with spacing $1/M$, both reduce to standard
softmax attention.

For \emph{operator self-attention},
$(\bm{Q}, \bm{K}, \bm{V})$ are linear projections of the same input observations, and
$\bm{\Delta\tau}$ are the quadrature weights of the observation's time grid. For \emph{operator cross-attention}, $\bm{Q}$ is projected from the query observations,
while $(\bm{K}, \bm{V})$ and $\bm{\Delta\tau}$ are projected from, and computed over,
the key-value observations and their corresponding time grids. Explicitly:
\begin{align}
\mathtt{SelfOpAttn}(\bm{x},\bm{\Delta \tau}, \bm{M}) &= {\tt OpAttn}(\bm{W_Q}\cdot \bm{x}, \bm{W_K}\cdot \bm{x}, \bm{W_V}\cdot \bm{x}\,; \bm{M}, \bm{\Delta\tau})\,,\\
\mathtt{CrossOpAttn}(\bm{x}, \bm{y},\bm{\Delta \tau}, \bm{M}) &= {\tt OpAttn}(\bm{W_Q}\cdot \bm{x}, \bm{W_K}\cdot \bm{y}, \bm{W_V}\cdot \bm{y}\,; \bm{M},   \bm{\Delta\tau})\,,
\end{align}
where $\bm{W_{ Q,K,V}}$ are the corresponding (learnable) linear transformations to the matrices $(\bm{Q}, \bm{K}, \bm{V})$, respectively. 


% ---------------------------------------------------------------------
\subsection{Flow-Time Conditioning and Triangular Past--Future Masking}
\label{app:flow_time}

\paragraph{Flow-time embedding.}
The ODE time $t \in [0,1]$ is encoded as a $d$-dimensional sinusoidal Fourier
feature vector with log-spaced frequencies,
\begin{equation}
    t_{\rm emb}
    = \Bigl[
        \sin(t\,\omega_1),\ldots,\sin(t\,\omega_{d/2}),\;
        \cos(t\,\omega_1),\ldots,\cos(t\,\omega_{d/2})
    \Bigr]
    \;\in\; \mathbb{R}^d,
    \label{eq:time_emb}
\end{equation}
where $\omega_k = f_{\max}^{-(k-1)/(d/2-1)}$ for $k = 1,\ldots, d/2$,
giving frequencies log-spaced between 1 and $1/f_{\max}$.
The embedding $t_{\rm emb}$ is added residually to every feature embedding in both encoder
and decoder after each attention block, ensuring the velocity field is a
smooth function of the ODE time at every network layer.
Note that $t$ is the \emph{flow} ODE time and plays no role in the observation
time axis $\tau$; the two are kept strictly separate throughout.

\paragraph{Reference measure and GP posterior.}
\label{app:ref_measure}
As formalised in Eq.~(14) of the main text, the reference measure
$\eta_F(\cdot \mid o, p)$ depends on whether a past prefix is available.
\begin{itemize}
\item In \emph{population synthesis} ($p = 0$), the flow's source is sampled from a zero-mean GP prior $\mathbf{z}_0 \sim \mathrm{GP}(0,\, k)|_{\boldsymbol{\tau}}$ with the standard RBF kernel evaluated at the target's time grid:
\begin{equation}\label{eq:rbf_kernel}
k(\tau,\tau') = \nu^2 \exp \bigl(-\frac{(\tau-\tau')^2}{4\ell^2}\bigr)\,.
\end{equation}
Here, the variance $\nu^2$ and the length scale $\ell$ are model hyperparameters. Since the (normalized) concentrations take values in $\mathbb{R}^+$ we apply a softplus transformation $\mathbf{z}_0 \to {\tt \log(1+\exp(\mathbf{z}_0 ))}$ to the output of the GP that guarantees positive concentrations values along the flow trajectories. 

\item In \emph{forecasting} ($p > 0$), the flow's source $\mathbf{z}_0 = (\mathbf{y}_{\tt P}, \mathbf{x}_{\tt F})$ is sampled from the GP
\emph{posterior} conditioned on the observed past $(\boldsymbol{\tau}_{\tt P}, \mathbf{y}_{\tt P})$:
\begin{equation}
    \mathbf{x}_{\tt F} \sim \mathrm{GP}\bigl(\bm{\mu}_*(\boldsymbol{\tau}_{\tt F}),\;
    \bm{\Sigma}_*(\boldsymbol{\tau}_{\tt F}, \boldsymbol{\tau}_{\tt F})\bigr),
    \label{eq:gp_post}
\end{equation}
where $\bm{\mu}_*$ and $\bm{\Sigma}_*$ are the posterior mean and covariance obtained via GP regression with \eqref{eq:rbf_kernel}. Furthermore, we apply a softplus transformation to the future $\bf x_{\tt F}$, leaving the past unchanged. The posterior already encodes continuity and monotonic behaviour compatible with early-phase PK profiles, providing a warm start for the flow that
substantially reduces the variance of the flow-matching objective compared to
an uninformed prior.
In both cases, past-slot values are fixed to the observed concentrations in
both $\mathbf{z}_0 = (\mathbf{y}_{\tt P}, \mathbf{x}_{\tt F})$ and $\mathbf{z}_1 = (\mathbf{y}_{\tt P}, \mathbf{y}_{\tt F})$
(see Eq.\ (15)), so the interpolation is trivial for past observations.
\end{itemize}

\paragraph{Triangular past--future masking.}
The prefix-dependent mask $\mathbf{M}_p$ in the flow-matching loss Eq.~(16) zeros out the velocity field at all past observations:
\begin{equation}
    v^\theta(t, \mathbf{z}_t, \mathcal{S})
    \;\leftarrow\;
    v^\theta(t, \mathbf{z}_t, \mathcal{S}) \odot \mathbf{M}_p,
    \qquad
    \mathbf{M}_p = \mathbf{1}\!\bigl[\tau > \tau_p\bigr].
    \label{eq:tri_mask}
\end{equation}
This is the discrete implementation of the triangular map structure from
Definition~3.1: no probability mass is transported at past time points, and
their values therefore remain exactly equal to the observed prefix
irrespective of the ODE time $t$.
Combined with the GP-posterior source~\eqref{eq:gp_post}, the flow
learns to move probability mass only in the unobserved future, while the past
acts as a fixed conditioning signal entering through the decoder transformer
~\eqref{app:decoder}.

\subsection{Context Study encoder}

A study $\mathcal{S} = \{\mathcal{D}^1, \ldots, \mathcal{D}^I\}$ contains
$I$ individuals. After normalization (described below), individual $i$ contributes
$C$ time slots (zero-padded where fewer observations exist), yielding a total number of observations $N = I \cdot C$. Each observation in the context study is described by $\mathcal{D}_j^i =({\tau}^i_j,\;{y}^i_j,\;a^i,\;r^i\bigr) \in \mathbb{R}^{4}$, $i\in[I]$, $j\in[C]$, where ${\tau}^i_j$ and ${y}^i_j$ are observation time and concentration, and $(a^i, r^i)$ are the dose amount and administration route. Concentrations are normalized per-study by the maximum concentration observed across all context individuals and time points; times are normalized by the maximum observation time in the study. Both scale factors are estimated from the context set $\mathcal{S}$ and applied consistently to all context and target quantities. The varying number of observation times per individual are tracked by a binary mask $\bm{m}^i\in \{0,1\}^C$ distinguishing true observations from zero-padded entries.

\paragraph{Input embedding}
Each study observation $\mathcal{D}_j^i$ is lifted to the
model dimension via a three-layer MLP $\Phi_{\mathcal S}:\mathbb{R}^4\to \mathbb{R}^d$ interleaved with $\rm GELU$ activations and combined with
a flow-time embedding (Section~\ref{app:flow_time}):
\begin{equation}
    \bm{h}^{i\,(0)}_j
    = \mathtt{LN} \Bigl[
        \Phi_{\mathcal{S}}\!\left(\mathcal{D}_j^i\right)
        + t_{\rm emb}
    \Bigr]
    \;\in\; \mathbb{R}^d,
    \label{eq:enc_embed}
\end{equation}
where $\tt LN(\cdot)$ stands for layer norm.

\paragraph{Encoder}
The study encoder applies $L_E$ stacked operator self-attention blocks $\Psi_\mathcal{S}^{\ell}:\bm{h}^{(\ell-1)}\to\bm{h}^{(\ell)}$, where we denote by ${\bm h}^{(\ell)} \in \mathbb{R}^{N \times d}$ the hidden state after block $\ell\in L$:
\begin{align}
    \bm h &\;=\; \bm{h}^{(\ell-1)}
    + \mathtt{SelfOpAttn}\bigl(\mathtt{LN}(\bm{h}^{(\ell-1)}),\;\bm{\Delta\tau},\;\bm{M}\bigr),
    \label{eq:enc_sa}\\
    \bm{h}^{(\ell)} &\;=\; \bm{h} + \Phi^{\ell}\bigl({\tt LN}(\bm{h})\bigr) + t_{\rm emb}\,\bm{1}^\top_N\,.
    \label{eq:enc_ffn}
\end{align}
Here, $\Phi^{\ell}$ is a two-layer GELU MLP with expansion ratio $4\times$ and $\bm\Delta\tau$ denotes the subject-specific time increments used to form the trapezoidal quadrature weights (Eq.~\eqref{eq:trap_weights}). To make the study-encoding depend explicitly on the functional flow-matching time $t$, we add after each block the time embedding $t_{\rm emb}$ broadcast to each observation in the study. Since subjects in the study are exchangeable but observations between subjects should not mix in the encoder, we use an attention mask $\bm M\in\mathbb{R}^{N\times N}$ that enforces \emph{subject-wise} attention only: we apply a block-diagonal pattern so that tokens from subject \(j\) may attend only to tokens from the same subject \(j\), preventing information flow across subjects within the encoder self-attention. Finally, the output of the encoder yields the hidden study representation:
\begin{equation}
\bm h^{\mathcal S} = {\tt LN}\bigl[\Psi^{L_E-1}_\mathcal{S}(\bm h^{(L_E-1)})\,\circ\,\cdots\circ\Psi^0_\mathcal{S}(\bm h^{(0)})\bigr]\,.
\end{equation}

\subsection{Target Individual Decoder}
\label{app:decoder}
For the observations of a new individual $\mathcal{D}_j = (\tau_j, y_j, a, r)\in \mathbb{R}^4$,
the decoder operates over an irregular grid of $T$ time points formed by concatenating the
observed prefix times $\boldsymbol{\tau}_{\tt P}$ and the future query
times $\boldsymbol{\tau}_{\tt F}$, sorted chronologically.
At flow time $t$, the flow interpolated concentration can be written as
\begin{equation}
    \mathbf{z}_t \;= \;({\bf y}_{\tt P}, {\bf y}_{t,\tt F}) \;= \;\bigl[t\,({\bf y}_{\tt P}, {\bf y}_{\tt F})\; +\; (1-t)\,({\bf y}_{\tt P}, \bf x_{\tt F})\; + \;\sigma_{\min}\,\bm \varepsilon\bigr] \odot \mathbf{M}_p
    \label{eq:zt_def}
\end{equation}
where $\bm\varepsilon \sim \mathcal{N}(0,\bm 1)$ and $\sigma_{\rm min}$ is a small hyperparameter that regularizes the conditional Dirac paths. The irregularity of the observation time grid for target individuals is kept track by a binary mask $\bm{m}\in \{0,1\}^T$ distinguishing true observations from zero-padded entries.

\paragraph{Input embedding.}
Each target individual observation $\mathcal{D}_j \in \mathbb{R}^4$ is embedded by a separate three-layer GELU MLP interleaved with GELU activations and added to the flow-time embedding:
\begin{equation}
    \bm{g}^{(0)}_j
    = \mathtt{LN} \Bigl[
        \Phi_{\mathcal{T}}\!\left(\mathcal{D}_j\right)
        \;+\; t_{\rm emb}
    \Bigr]
    \;\in\; \mathbb{R}^d,
    \label{eq:enc_embed}
\end{equation}
where $\tt LN(\cdot)$ stands for layer norm.

\paragraph{Decoder}
The decoder applies $L_D$ blocks, each containing two attention components with non-causal structure: an intra-individual operator self-attention layer followed by a operator cross-attention layer over the context study $\mathcal{S}$. Writing $\bm{g}^{(\ell)} \in \mathbb{R}^{T \times d}$ for the hidden state after block $\ell\in L_{D}$, the decoder block $\Psi_\mathcal{T}^{\ell}:\bm{g}^{(\ell-1)}\to\bm{g}^{(\ell)}$ is defined by
\begin{align}
    \bm g &\;=\; \bm{g}^{(\ell-1)}
   \; +\; \mathtt{SelfOpAttn}\bigl(\mathtt{LN}(\bm{g}^{(\ell-1)}),\;(\bm{\Delta\tau}_{\tt P}, \bm{\Delta\tau}_{\tt F}),\;\bm{M}_1\bigr),
    \label{eq:enc_sa}\\
    \bm g^\prime &\;=\; \bm{g} \;+\; \Phi_1^{\ell}\bigl({\tt LN}(\bm{g})\bigr)\,,\\
    \bm g^{\prime\prime} &\;=\; \bm{g^{\prime}}
    \;+\; \mathtt{CrossOpAttn}\bigl(\mathtt{LN}(\bm g^\prime)\,,\; \bm{h}^{\mathcal{S}},\;\bm{\Delta\tau},\;\bm{M}_2\bigr)\\
    \bm g^{(\ell)} &\;=\; \bm{g}^{\prime\prime} \;+ \;\Phi_2^{\ell}\bigl({\tt LN}(\bm{g}^{\prime\prime})\bigr)\; + \;t_{\rm emb}\,\bm{1}^\top_T,
\end{align}

where $(\bm\Delta\bm\tau_{\tt P},\bm\Delta\bm\tau_{\tt F})$ and $\bm\Delta\bm\tau$ are the irregular times increments of the (past, future) target individual and study context, respectively, $\Phi_{1,2}$ are 2-layered MLPs and $\bm{h}^\mathcal{S}$ is the study encoder output used for the cross-attention key and values. The attention mask matrices, $\bm M_1\in {\mathbb{R}^{T\times T}}$ and $\bm M_2\in {\mathbb{R}^{T\times N}}$, are defined through their respective binary observation masks, preventing queries to attend to invalid zero-padded entries. This setup allows for information to flow freely between the target's future observations $\bf{y}_{\tt F}$ and the past contextual observations $\bf{y}_{\tt P}$ during self-attention, while subsequently cross-attending to all study observations at all possible times. The output of the target individual decoder reads:
\begin{equation}
\bm g^{\mathcal T} = {\tt LN}\bigl[\Psi^{L_D-1}_\mathcal{T}(\bm g^{(L_D-1)})\,\circ\,\cdots\circ\Psi^0_\mathcal{T}(\bm g^{(0)})\bigr]\,.
\end{equation}

\paragraph{Output head.} Finally, the decoder's output feeds into the head of the architecture parametrized by a 3-layered MLP $\Phi_{\rm head}$:
\begin{equation}
    \bigl[v^\theta_t(\mathbf{z}_t,\mathcal{S})\bigr]_j
    = \mathrm{\Phi}_{\mathrm{head}}\bigl(\bm g^{\mathcal{T}}_j\bigr)\,.
    \label{eq:head}
\end{equation}
This yields the functional velocity field evaluated at all query times.


% ---------------------------------------------------------------------
\subsection{B.7 \quad Hyperparameters and Training Details}
\label{app:hyperparams}

Table~\ref{tab:hyperparams} summarizes the architectural and training
hyperparameters used for the \texttt{PFF} results reported in Section~5.

\begin{table}[h]
\centering
\caption{Hyperparameters of \texttt{PFF}.}
\label{tab:hyperparams}
\begin{tabular}{lc}
\toprule
\textbf{Parameter} & \textbf{Value} \\
\midrule
\multicolumn{2}{l}{\textit{Architecture}} \\
Hidden dimension $d$                        & 256 \\
Encoder depth $L_E$                         & 4   \\
Decoder depth $L_D$                         & 4   \\
Attention heads (encoder / decoder) $H$     & 4 / 4 \\
FFN expansion ratio                         & $4\times$ \\
Dropout                                     & 0.1 \\
Fourier max.\ frequency $f_{\max}$          & 256 \\
\midrule
\multicolumn{2}{l}{\textit{Reference measure (GP regression)}} \\
Kernel variance $\nu^2$                     & $1\!\times\!10^{-4}$ \\
Length scale $\ell$                         & $1.7\!\times\!10^{-3}$ \\
Cholesky Jitter $\epsilon$                           & $1\!\times\!10^{-7}$ \\
Output transform                            & Softplus \\
\midrule
\multicolumn{2}{l}{\textit{Flow matching}} \\
Path noise $\sigma_{\min}$                  & $10^{-4}$ \\
ODE integration steps (inference)           & 100 \\
\midrule
\multicolumn{2}{l}{\textit{Optimisation}} \\
Optimizer                                   & AdamW \\
$(\beta_1,\,\beta_2)$                       & $(0.9,\;0.999)$ \\
Learning rate                               & $1\!\times\!10^{-5}$ \\
Gradient clip (global $\ell_2$ norm)        & 0.5 \\
Batch size                                  & 64 \\
Epochs                                      & 300 \\
Warm-up epochs                              & 5 \\
\bottomrule
\end{tabular}
\end{table}







% We now introduce the neural network architecture details for the Amortized In Context Mixed Effect Transformer (AICMET) model, required to define both encoder \eqref{eq:encoder} and decoder \eqref{eq:decoder}. We can see that in order to create a faithful representation of the data, we first need to create a longitudinal representation per patient that respects the time information per individual. 
% Then, we aggregate the time dimension to obtain a representation per individual, and finally we aggregate all the individuals to obtain one representation per study. 
% In the following, we use  \(H\) as hidden dimension and \(Z_{d}\) as latent dimension for the unstructured latent variables. A superscript in parentheses, e.g.\ \(\mathbf{h}^{(1)}\), denotes the
% \emph{layer index}.



% \vspace{4pt}
% \noindent\textbf{} We denote by $\phi^{\theta}(\cdot)$ a linear projection $\mathbb{R}^{m} \to \mathbb{R}^{H}$, by $\Phi^{\theta}(\cdot)$ a feed-forward network $\mathbb{R}^{H} \to \mathbb{R}^{H}$, by $\rho^{\theta}(\cdot)$ a recurrent layer (GRU/LSTM), by $\psi^{\theta}(\mathbf{q},\mathbf{K},\mathbf{V})$ a multi-head attention module, and by $\Psi^{\theta}(\mathbf{Q},\mathbf{K},\mathbf{V})$ a transformer encoder (linear attention).
% For both $\psi^\theta(\cdot, \cdot, \cdot)$ and $\Psi^\theta(\cdot, \cdot, \cdot)$ denote transformer encoders with linear attention \citep{katharopoulos2020transformers}, both of which take three arguments as inputs (i.e., queries, keys and values). 
% We use the dot-product softmax kernel, i.e.,
% \(
% \psi^\theta(\mathbf{Q}, \mathbf{K}, \mathbf{V}) := \operatorname{softmax}\left( \frac{\mathbf{Q} \mathbf{K}^\top}{\sqrt{d_k}} \right) \mathbf{V}.
% \)
% Since PK data are often sparsely and irregularly sampled, we expand the features by defining the individual dataset as
% \(
% \widetilde{\mathcal{D}}^{i}
%   =\bigl\{(y_{j},\tau_{j},\Delta\tau_{j})\bigr\}_{j=1}^{T_{i}-1}$ wiht $\Delta\tau_{j}=\tau^{i}_{j+1}-\tau^{i}_{j}.
% \)
% This simplifies the task of the network of handling irregular time intervals, as this representation enforces this feature.




% We first introduce the encoder, which we construct by first introducing the individual representation and then aggregating for the studies representations. 
% %It consist of an embedding layer 
% \paragraph{Per-transition embeddings.}
% We pass the basic features trough a linear layer and concatenate, thereby defining a basic data embedding: 
% \[
% \mathbf{e}^{(0),i}_{j}
%  =\text{concat}\Bigl[
%      \phi^{\theta}_{y}(y_{j}),
%      \phi^{\theta}_{\tau}(\tau_{j}),
%      \phi^{\theta}_{\Delta\tau}(\Delta\tau_{j})
%    \Bigr]\in\mathbb{R}^{H}.
% \]
% \paragraph{Longitudinal Representation.} To capture the dynamical information from the embeddings, we now create a representation per individual and per time step, using either a recurrent neural network (GRU/LSTM) or a transformer architecture:
% \[
% \mathbf{h}^{(1),i}_{j}
% =\rho^{\theta}\bigl(\mathbf{e}^{(0),i}_{1:j}\bigr)\in\mathbb{R}^{H}.
% \]
%  Next, in order to judge the overall curve shape,  a self-attention mechanism is applied over the longitudinal representations $\mathbf{h}^{(1),i}_{j}$. 
%  This allows the different native features of the drug profiles, such as the peak and slope of different sections around the peak, to be weighted and compared against each other. We first stack the representations
% \(
% \mathbf{H}^{(1),i}=[\mathbf{h}^{(1),i}_{1},\dots,\mathbf{h}^{(1),i}_{T_{i}-1}]
%  \in\mathbb{R}^{H\times(T_{i}-1)}
% \); then
% \[
% \mathbf{H}^{(2),i}
%  =\Psi^{\theta}\!\bigl(\mathbf{H}^{(1),i},
%                        \mathbf{H}^{(1),i},
%                        \mathbf{H}^{(1),i}\bigr).
% \]
% \paragraph{Individual posterior.}
% Now we aggregate with attention pooling so that we are able to obtain a representation per individual
% \[
% \mathbf{c}_{i}
%  =\psi^{\theta}\!\bigl(\mathbf{q},
%                        \mathbf{H}^{(2),i},
%                        \mathbf{H}^{(2),i}\bigr)\in\mathbb{R}^{H}.
% \]
% We then differentiate the mean and variance of the encoder with a final MLP layer
% \[
% q(\bz_{i}\mid\mathcal{D}^{i})
%  =\mathcal{N}\!\Bigl(
%      \bz_{i}\in\mathbb{R}^{Z_{d}}
%      \;\bigl|\;
%      \Phi^{\theta}_{\mu}(\mathbf{c}_{i}),
%      \operatorname{diag}\!\bigl(
%        e^{\,\Phi^{\theta}_{\sigma}(\mathbf{c}_{i})}
%      \bigr)
%    \Bigr).
% \]
% \paragraph{Study Posterior.}
% To obtain a study-specific representation, we concatenate the individual summaries
% \(
% \mathbf{C}=[\mathbf{c}_{1},\dots,\mathbf{c}_{I}]
%  \in\mathbb{R}^{H\times I},
% \)
% and then summarize with another attention pooling mechanism,
% \(
% \mathbf{c}_{s}
%   = \psi^{\theta}\bigl(\mathbf{q}_{s}, \mathbf{C}, \mathbf{C}\bigr).
%   \)
% Finally the study encoder is obtained by
% \[
% q(\bz_{s} \mid \mathcal{S})
%   = \mathcal{N}\!\Bigl(
%      \bz_{s} \in \mathbb{R}^{Z_{d}} \;\bigl|\;
%      \Phi^{\theta}_{\mu}(\mathbf{c}_{s}),
%      \operatorname{diag}\!\bigl(
%        e^{\,\Phi^{\theta}_{\sigma}(\mathbf{c}_{s})}
%      \bigr)
%    \Bigr).
% \]

% \subsection{Decoder}
% \vspace{-0.08cm}
% To design a decoder suited for PK data, we must accommodate the irregular sampling schedules typically observed in clinical datasets while simultaneously exploiting the strong \emph{context-learning} capabilities of modern transformers. Conventional neural ODE approaches \citep[e.g.][]{lu2021deep} rely on adjoint-sensitivity computations; this can hinder scalability for a data prior like ours that requires many samples to enforce expert based inductive bias. 
% We therefore adopt a transformer decoder endowed with \emph{functional attention} \citep{seifner2025foundation}. The requested prediction/sample time points $\tau$ are embedded as \textbf{queries}, whereas the dosing information $\mathbf{u}^n$ together with the population-level ($\mathbf{z}_{s}$) and individual-level ($\mathbf{z}_{n}$) latent variables are embedded as \textbf{keys} and \textbf{values}. Through self-attention, the model acts as a context learner: each query time dynamically attends to the most informative dosing and latent-effect context, thereby defining a distribution over concentration-time functions conditioned on both, $\mathbf{z}_{s}$ and $\mathbf{z}_{i}$. For a new subject $\mathcal{D}^{(n)}$, we first sample $\mathbf{z}_{n}$ from its prior, then propagate the query-time set through the decoder to obtain the predictive concentration distribution at each requested time point.
% Specifically, for a new individual \(\mathcal{D}^{n}\), we first sample
% \(
% \bz_{n}\sim q(\bz_{n}|\mathcal{D}^{n})\) and \(\bz_{s}\sim q(\bz_{s}|\mathcal{S})
% \) for prediction, or \(\bz_n \sim \mathcal{N}(0,I)\) for generation
% and embed a query time \(\tau\): 
% \begin{align}
% \mathbf{q}_{\tau} 
%   &= \phi^{\theta}_{\text{dec}}(\tau) \in \mathbb{R}^{H}, \nonumber \\
% \mathbf{K}_{n} 
%   &= \Phi^{\theta}_{K}\left([\bz_{n}; \bz_{s}; \mathbf{u}]\right) \in \mathbb{R}^{H \times 1}, \nonumber  \\
% \mathbf{V}_{n} 
%   &= \Phi^{\theta}_{V}\left([\bz_{n}; \bz_{s}; \mathbf{u}]\right) \in \mathbb{R}^{H \times 1}, \nonumber  \\
% \mathbf{c}_{\tau}
%   &= \psi^{\theta}\left(\mathbf{q}_{\tau}, \mathbf{K}_{n}, \mathbf{V}_{n}\right) \in \mathbb{R}^{H}.
% \end{align}
% It only remains to define the final mean and variance heads,
% \(
% (\mu_{\tau},\log\sigma^{2}_{\tau})
%  =\Phi^{\theta}_{\text{dec}}(\mathbf{c}_{\tau}),
% \) and specify the distribution to a diagonal Gaussian 
% \(
% p\bigl(y_{\tau}\mid\tau,\bz_{n},\bz_{s}\bigr)
% =\mathcal{N}\bigl(y_{\tau}\mid\mu_{\tau},\sigma^{2}_{\tau}\bigr).
% \)
